\section*{Resumen}

La arquitectura en capas ha sido uno de los enfoques estructurales más influyentes en el desarrollo de software, lo cual se debe a la claridad que aporta al separar responsabilidades, facilitar la mantenibilidad y promover la modularidad. Sin embargo, a lo largo del tiempo, múltiples investigaciones han mostrado tanto la relevancia histórica del patrón como sus límites frente a los sistemas distribuidos, móviles, orientados a datos y basados en aprendizaje automático. Este artículo examina el enfoque en capas  y revisa los desafíos técnicos que enfrenta en escenarios actuales.

En este artículo, se identificaron estudios procedentes de investigaciones recientes que incluyen técnicas de reconstrucción arquitectónica, modelos de visualización estructural, propuestas de arquitectura limpia, análisis del rendimiento mediante modelos probabilísticos y evaluaciones empíricas en sistemas modernos. A partir de esto, se identifica cómo la erosión arquitectónica, la latencia del flujo secuencial y las violaciones de dependencia representan riesgos persistentes para la sostenibilidad del patrón.

Además, se presenta un análisis aplicado a un sistema real: una plataforma de carnetización digital que integra autenticación, control de asistencia, cargas masivas de información y notificaciones automatizadas, donde se explora cómo la arquitectura por capas puede servir como fundamento estructural y también se identifican los ajustes necesarios para mantener su rendimiento y su capacidad de evolución en un ecosistema basado en procesamiento intensivo de datos y flujos operativos concurrentes.

Los resultados permiten comprender el papel actual del enfoque por capas en la ingeniería de software, destacando tanto su relevancia como los retos que deben atenderse para garantizar su eficacia en sistemas modernos.
\textbf{Palabras clave:} arquitectura por capas, mantenibilidad, rendimiento, refactorización, aprendizaje automático, modularidad.
